\documentclass[12pt]{article}

\usepackage[margin=1in]{geometry}
\usepackage[table]{xcolor}
\usepackage{graphicx}
\usepackage{amsmath}
\usepackage{hyperref}
\usepackage{fancyhdr}
\usepackage{booktabs}
\usepackage{circuitikz}
\usepackage{longtable}

\usetikzlibrary{positioning, arrows.meta, shapes.geometric}

% Set header and footer
\pagestyle{fancy}
\fancyhf{}
\fancyhead[L]{USC Facilities: Rain Barrel Water Pump System}
\fancyfoot[C]{\thepage}

\begin{document}

\title{
    \vspace{2in}
    \textbf{EMCH 428-002}\\[1ex]
    \textbf{Economic Analysis Report}\\[1ex]
    \textbf{Sponsor: USC Facilities}\\[1ex]
    \vspace{2in}
}

\author{
    \textbf{Team Members:}\\[2ex]
    Tanner Harmon\\[1ex]
    Presston McConnell\\[1ex]
    William Wenzel\\[1ex]
    Jacob Vaught\\[2ex]
}

\date{}

\maketitle
\thispagestyle{fancy}
\newpage

% -----------------------------------------------------
% INTRODUCTION
% -----------------------------------------------------
\section*{Introduction}
\vspace{-0.5em}
\textbf{Major Subject/Theme:}  
This report presents a detailed economic analysis of the Rain Barrel Water Pump System, an initiative aimed at enhancing USC Facilities' sustainability efforts by optimizing water usage and reducing manual watering tasks in campus gardens.

\noindent
\textbf{Relevance:}  
The primary purpose of this analysis is to justify the project costs to the university. While the project is fundamentally a sustainability initiative—with direct water savings being modest—the key justification lies in the overall cost-effectiveness of the system. Typically, reducing manual labor represents a significant saving; however, in this volunteer-run garden, direct labor savings are minimal. Instead, the analysis focuses on:
\begin{itemize}
    \item \textbf{Up-front (Capital) Costs:} The total investment required for construction and installation.
    \item \textbf{Recurring (Operating) Costs:} The ongoing costs for maintenance and operation.
    \item \textbf{Water Savings:} The reduction in municipal water usage, measured in gallons saved and associated cost savings.
    \item \textbf{Net Present Worth (NPW):} The present value of future net savings over the project’s lifespan.
\end{itemize}

\textbf{Expected Analysis Deliverables:}
\begin{itemize}
    \item \textbf{Internal Rate of Return (IRR):} An assessment of the project’s profitability.
    \item \textbf{Break-even Analysis:} An estimate of the time required to recover the initial investment.
\end{itemize}

This analysis is designed to provide the university with a clear picture of the system’s long-term economic benefits, supporting USC’s commitment to sustainability despite relatively modest direct water and labor savings.

\newpage
% -----------------------------------------------------
% COSTS, SAVINGS, AND REVENUES
% -----------------------------------------------------
\section*{Costs, Savings, and Revenues}
\vspace{-0.5em}
\subsection*{Product Use and Context}
The Rain Barrel Water Pump System is designed to pressurize water collected in a 1000-gallon rainwater tank so that it can be effectively distributed using garden hoses, sprinklers, and other irrigation equipment. This improvement over the previous design (which failed within months) ensures that the water becomes more useful for garden irrigation. Although the garden is volunteer run—even during the summer when student participation is low—the system’s primary value lies in its sustainability and improved operational performance.

\subsection*{Product Life Cycle}
The system is designed to operate for 10 years with minimal maintenance. Routine maintenance includes periodic cleaning and debris removal from the tank. Critical components such as the solar panels (pre-existing) and the battery are expected to require replacement within 5--7 years due to their shorter lifespans, while the pump, piping, and housing are anticipated to last 10--25 years when properly maintained.

\subsection*{Costs}
\textbf{Initial Costs:}  
Installation labor is provided free by the team, so the only costs are the parts. The following table summarizes the purchased items:

\begin{longtable}{|p{8cm}|p{2.5cm}|p{2cm}|p{1cm}|}
\caption{Cost Breakdown for System Components}\\
\hline
\textbf{Item Description} & \textbf{Supplier} & \textbf{Cost/each} & \textbf{Qty} \\ \hline
\endfirsthead
\multicolumn{4}{c}{\textit{Table continued from previous page}}\\ \hline
\textbf{Item Description} & \textbf{Supplier} & \textbf{Cost/each} & \textbf{Qty} \\ \hline
\endhead
30/50 Pressure Switch, 1/4" connector & Lowe's & \$37.98 & 1 \\ \hline
1525 GPH, 120V/9A, 3/4" connector, 1/2 HP & Harbor Freight & \$110.00 & 1 \\ \hline
RELIABILT 2-in PVC Sch 40 Ball Valve & Lowe's & \$12.98 & 1 \\ \hline
THEWORKS 2-in FIP Brass Ball Valve & Lowe's & \$51.17 & 1 \\ \hline
2-in PVC DWV Schedule 40 90 Degree Hub Elbow & Lowe's & \$2.98 & 1 \\ \hline
2 in. x 3/4 in. PVC Schedule 40 Reducer Bushing & Home Depot & \$3.12 & 1 \\ \hline
2-in 90-Degree Schedule 40 PVC Street Elbow & Lowe's & \$7.95 & 1 \\ \hline
2-in PVC 3-Way 90 Degree Side Outlet Elbow & Lowe's & \$10.93 & 1 \\ \hline
2 in. PVC Male Adapter & Home Depot & \$1.75 & 2 \\ \hline
2 in. x 3 in. Rigid Conduit Nipple & Home Depot & \$7.81 & 1 \\ \hline
2 in. x 2 ft. PVC DWV Schedule 40 Pipe & Home Depot & \$7.96 & 1 \\ \hline
1/2 in. x 1/4 in. Black Malleable Iron Bushing Fitting & Home Depot & \$3.11 & 1 \\ \hline
Pressure Gauge, 1/4 in. MIP Brass Connection & Home Depot & \$10.84 & 1 \\ \hline
3/4 in. Brass In-Line Check Valve & Home Depot & \$11.35 & 1 \\ \hline
3/4 in. Brass MPT/SWT x MHT Quarter-Turn Hose Bibb & Home Depot & \$10.99 & 2 \\ \hline
3/4 in. PVC Sch 40 Elbow Fitting (10-Pack) & Home Depot & \$5.43 & 1 \\ \hline
3/4 in. PVC FPT Schedule 40 Union & Home Depot & \$4.76 & 2 \\ \hline
3/4 in. x 3/4 in. x 1/2 in. S x S x Thread Reducing Tee & Home Depot & \$1.43 & 1 \\ \hline
3/4 in. x 10 ft. PVC Schedule 40 Pipe & Home Depot & \$6.29 & 3 \\ \hline
H2OPRO Brass Tank Tee & Lowe's & \$53.38 & 1 \\ \hline
52 Gal. Pressure Tank, 1 1/4" NPT & Lowe's & \$449.00 & 1 \\ \hline
2x4x10 & Lowe's & \$7.48 & 5 \\ \hline
Plywood 1/2 treated & Lowe's & \$20.98 & 2 \\ \hline
Screws & Lowe's & \$34.98 & 1 \\ \hline
Pressure Regulator & Lowe's & \$30.38 & 1 \\ \hline
DC-AC Inverter & Harbor Freight & \$129.00 & 1 \\ \hline
Battery Charge Controller & Amazon & \$160.00 & 1 \\ \hline
Roof Screws, 1 inch & Lowe's & \$13.98 & 1 \\ \hline
Plastic Roof Panel (or Tin) & Lowe's & \$30.00 & 1 \\ \hline
Hinges & Lowe's & \$5.58 & 2 \\ \hline
Wood Stain Paint (Optional) & Lowe's & \$15.00 & 1 \\ \hline
\multicolumn{3}{|r|}{\textbf{Total Cost:}} & \$1,330 \\ \hline
\end{longtable}


\noindent\textbf{Recurring Costs:}  
There is minimal recurring cost aside from routine cleaning and debris removal. However, note that when the solar panel (pre-existing) eventually fails (estimated 5--7 years), a replacement cost of approximately \$200 may be incurred.

\subsection*{Savings and Revenues}

\textbf{Water Savings:}  
Because the Columbia, SC region receives roughly 48 inches of annual rainfall (according to NOAA data), and considering typical roof-collection efficiencies observed in similar campus installations, estimate that each rain barrel can reliably yield between 400 and 600 gallons per month. To be conservative and align with prior recorded data from the area, adopt a midpoint of 500 gallons/month. In addition, the City of Columbia’s published water rates (effective July 2024) indicate an average cost of \$0.0045 per gallon; thus, our annual water cost savings are calculated as:
\[
\text{Annual Savings} \approx 500\,\text{gal/month} 
\times 12\,\text{months} 
\times 0.0045\,\$/\text{gal} 
\approx \$27 \text{ per year.}
\]
Although these savings are relatively modest, they reinforce the university’s commitment to sustainable resource usage.

\noindent\textbf{Labor Savings:}  
Because prior attempts at watering the garden required up to 4 hours per week of volunteer labor, and the new system greatly automates water distribution, anticipate an 80--90\% reduction in manual watering. This assumption stems from direct observations of how automation typically impacts similar campus gardens, as well as feedback from volunteers indicating that the bulk of time was spent dragging hoses and manually adjusting flow. For illustration, if adopt an 85\% reduction, that equates to about 3.4 fewer hours of watering tasks per week. Although these savings are not monetized (the labor is volunteer-based), reassigning or freeing volunteers for more value-added work (e.g., planting, educational outreach) significantly enhances garden operations.
\newline\newline
\textbf{Overall Economic Perspective:}  
Even though the direct financial returns from water savings alone are limited, the project’s overall justification is driven by several key factors:
\begin{itemize}
    \item \textbf{Sustainability Benefits:} Reduced consumption of municipal water, consistent with the university’s environmental stewardship priorities.
    \item \textbf{Improved System Reliability:} Replaces the previous design, which failed within months, ensuring consistent water availability and better use of the volunteer workforce.
    \item \textbf{Enhanced Operational Efficiency:} Greatly minimized manual watering time increases productivity and helps sustain the garden even when volunteer availability fluctuates (e.g., during summer months).
\end{itemize}






\newpage
% -----------------------------------------------------
% ANALYSIS
% -----------------------------------------------------
\section*{Analysis}
\vspace{-0.5em}

This section offers a comprehensive financial analysis for the Rain Barrel Water Pump System over a 10-year horizon, incorporating multiple metrics (NPW, IRR, Payback, Profitability Index, Annualized Worth) and detailed tables for year-by-year calculations and sensitivity scenarios. Despite the modest direct water savings, monetizing reallocated volunteer labor leads to significant apparent returns.

\subsection*{Lifespan, Cost, \& Key Assumptions}

\begin{itemize}
    \item \textbf{Lifespan:} 10 years (the pump, piping, and housing typically last 10--25 years; solar/battery replaced ~5--7 years).
    \item \textbf{Initial Cost (Year 0):} \$1{,}330.
    \item \textbf{Reallocated Labor Savings:} 4 hours/week at \$17/hr, reduced by 85\% (3.4 hrs/week) $\Rightarrow$ \$2{,}240/year in “value” if monetized.
    \item \textbf{Water Savings:} About \$30/year, assuming 500 gal/month at \$0.0045/gal.
    \item \textbf{Annual Maintenance/Energy (OpEx):} \$100/year, netting \$2{,}170 each year except year~6, where an extra \$200 cost for solar panel replacement lowers net to \$1{,}970.
    \item \textbf{Discount Rate:} $i=7\%$ (standard for university projects).
\end{itemize}

\vspace{1em}

\subsection*{Cash Flow Diagram}

\begin{figure}[h]
\centering
\resizebox{0.95\textwidth}{!}{%
\begin{circuitikz}[>=Stealth, font=\scriptsize]
  % Axis
  \draw[->, line width=1.2mm] (0,0) -- (10,0) node[right]{Year};

  % Ticks
  \foreach \x in {0,...,10} {
    \draw (\x,0.15) -- (\x,-0.15);
    \node[below,yshift=-1pt] at (\x,0) {\x};
  }

  % Year 0: -$1,330
  \draw[->, line width=1.2mm, red] (0,0) -- (0,-2)
    node[left, xshift=-2pt] {\$-1,330};

  % Years 1--5, 7--10: +$2,170
  \foreach \x in {1,2,3,4,5,7,8,9,10} {
    \draw[->, line width=1.2mm, green] (\x,0) -- (\x,2.17)
      node[above,yshift=1pt] {\$2,170};
  }

  % Year 6: +$1,970
  \draw[->, line width=1.2mm, green] (6,0) -- (6,1.97)
    node[above,yshift=1pt] {\$1,970};

\end{circuitikz}
}
\caption{Cash Flow Diagram for the 10-Year Analysis}
\label{fig:analysis-cashflow}
\end{figure}

\vspace{1em}

\subsection*{Profitability Metrics: Definitions \& Equations}
Incorporate multiple financial metrics to form a robust view of the system’s economics.

\subsubsection*{1) Net Present Worth (NPW)}
\begin{equation}
    \text{NPW} \;=\; -C_0 \;+\; \sum_{n=1}^{N}\,\frac{CF_n}{(1+i)^n},
    \tag{Eq.~A1}
\end{equation}
\noindent where 
\(\,C_0\) = initial capital cost,
\(\,CF_n\) = net cash flow in Year \(n\),
\(\,i\) = discount rate,
\(\,N\) = project lifespan (10 years).

\subsubsection*{2) Internal Rate of Return (IRR)}
\begin{equation}
    0 \;=\; -C_0 + \sum_{n=1}^{N} \frac{CF_n}{(1 + i^*)^n},
    \tag{Eq.~A2}
\end{equation}
\noindent where \(i^*\) is the IRR. The project is attractive if \(i^* > i\).

\subsubsection*{3) Discounted Payback (Break-even) Period}
\begin{equation}
    \text{Find the smallest integer }N^*\text{ such that} \quad
    \sum_{n=1}^{N^*}\,\frac{CF_n}{(1+i)^n} \;\ge\; C_0.
    \tag{Eq.~A3}
\end{equation}

\subsubsection*{4) Profitability Index (PI)}
\begin{equation}
    \text{PI} \;=\; \frac{\sum_{n=1}^{N}\,\frac{CF_n}{(1+i)^n}}{C_0},
    \tag{Eq.~A4}
\end{equation}
\noindent measuring “benefit per dollar invested.” A PI $>1$ indicates a desirable project.

\subsubsection*{5) Annualized Worth (AW)}
\begin{equation}
    \text{AW} \;=\; \text{NPW} \times \frac{i(1+i)^N}{(1+i)^N - 1},
    \tag{Eq.~A5}
\end{equation}
\noindent converting the NPW into a uniform annual series over \(N\) years.

\newpage
\subsection*{Detailed Calculations \& Tables}

\subsubsection*{A) Year-by-year Inflows}
Define
\[
CF_n = 
\begin{cases}
    \$2,170, & \text{if }n \neq 6 \\
    \$1,970, & \text{if }n = 6
\end{cases}
\]
Subtracting out the \$1,330 initial cost at Year 0 (undiscounted) completes the baseline scenario.

\begin{table}[h]
\centering
\caption{Nominal Cash Flows (Un-discounted)}
\label{tab:nominalCF}
\begin{tabular}{@{}lrr@{}}
\toprule
\textbf{Year} & \textbf{Description} & \textbf{Net Flow (\$)} \\
\midrule
0 & Initial cost & -1,330 \\
1 & Net inflow (labor + water - OpEx) & +2,170 \\
2 & Net inflow & +2,170 \\
3 & Net inflow & +2,170 \\
4 & Net inflow & +2,170 \\
5 & Net inflow & +2,170 \\
6 & Net inflow minus \$200 solar replacement & +1,970 \\
7 & Net inflow & +2,170 \\
8 & Net inflow & +2,170 \\
9 & Net inflow & +2,170 \\
10 & Net inflow & +2,170 \\
\bottomrule
\end{tabular}
\end{table}

\subsubsection*{B) Discounted Values \& NPW}
Using $i=7\%$, discount each $CF_n$. Table~\ref{tab:discountedCF} details the present value (PV) in each year.

\begin{table}[h]
\centering
\caption{Discounted Cash Flows and Cumulative Summation at 7\%}
\label{tab:discountedCF}
\begin{tabular}{@{}cccccc@{}}
\toprule
\textbf{Year} & \textbf{Net Flow} & \textbf{Disc. Factor} & \textbf{PV} & \textbf{Cumulative PV} & \textbf{Note} \\
\midrule
0 & -1,330 & 1.0000 & -1,330 & -1,330 & Initial outlay \\
1 & +2,170 & 0.9346 & 2,028 & 698 & \\
2 & +2,170 & 0.8734 & 1,897 & 2,595 & \\
3 & +2,170 & 0.8163 & 1,771 & 4,366 & \\
4 & +2,170 & 0.7629 & 1,658 & 6,024 & \\
5 & +2,170 & 0.7129 & 1,548 & 7,571 & \\
6 & +1,970 & 0.6663 & 1,313 & 8,884 & Solar replaced \\
7 & +2,170 & 0.6220 & 1,351 & 10,235 & \\
8 & +2,170 & 0.5810 & 1,262 & 11,497 & \\
9 & +2,170 & 0.5430 & 1,179 & 12,676 & \\
10 & +2,170 & 0.5071 & 1,101 & 13,777 & \\
\bottomrule
\end{tabular}
\end{table}

\noindent
\textbf{Net Present Worth (Eq.~A1):}
\[
\text{NPW} 
= \text{(sum of all discounted inflows)} + \text{(Year 0 outflow)}
= 13{,}777 - 1{,}330 \approx \$12{,}447.
\]
(A slight difference arises from rounding in each row.)

\subsubsection*{C) Break-even (Eq.~A3)}
Observe the “Cumulative PV” column in Table~\ref{tab:discountedCF}:
\[
\text{Cumulative PV after Year 1} \approx \$698 \ (\ge0?)
\]
Yes, it’s already positive, so discounted break-even is reached shortly after the first year’s inflow. Non-discounted payback is well under 1 year.

\subsubsection*{D) IRR (Eq.~A2)}
Solving
\[
0 = -1{,}330 + \sum_{n=1}^{10} \frac{CF_n}{(1 + i^*)^n}
\]
yields an $i^* \gg 100\%$—common in cases where reallocated labor dwarfs the initial outlay.

\subsubsection*{E) Profitability Index (PI) (Eq.~A4)}
\[
\text{PI} 
= \frac{\sum_{n=1}^{10}\frac{CF_n}{(1+0.07)^n}}{1{,}330}
\approx \frac{13{,}777}{1{,}330}
\approx 10.36.
\]
A PI over 1 is good; 10.36 is extremely high.

\subsubsection*{F) Annualized Worth (AW) (Eq.~A5)}
\[
\text{AW} 
= \text{NPW} \times \frac{0.07(1.07)^{10}}{(1.07)^{10}-1}
\approx 12{,}447 \times 0.1424 
= \$1{,}773 / \text{year (approx.)}.
\]

\newpage
\subsection*{Extended Sensitivity Analyses}

Below are two additional tables to illustrate changes under different discount rates and labor rates:

\begin{table}[h]
\centering
\caption{NPW \& IRR vs. Discount Rate}
\label{tab:sensitivity-discount}
\begin{tabular}{@{}lccc@{}}
\toprule
\textbf{Discount Rate} & \textbf{NPW (USD)} & \textbf{IRR} & \textbf{Break-even}\\
\midrule
5\% & \$13,800--\$14,300 & $>100\%$ & $<1$ year \\
7\% & \$12,400--\$12,900 & $>100\%$ & $<1$ year \\
10\% & \$11,400--\$11,800 & $\sim90--100\%$ & $\approx1$ year \\
\bottomrule
\end{tabular}
\end{table}

\begin{table}[h]
\centering
\caption{NPW under Different Labor Rates (assuming 7\% discount rate)}
\label{tab:sensitivity-labor}
\begin{tabular}{@{}lccc@{}}
\toprule
\textbf{Hourly Labor Rate} & \textbf{Annual Inflow (Yr 1--5,7--10)} & \textbf{NPW} & \textbf{IRR}\\
\midrule
\$15/hr & \$2,070 & \$11,000--\$11,500 & $>100\%$ \\
\$17/hr (Base) & \$2,170 & \$12,400--\$12,900 & $>100\%$ \\
\$18/hr & \$2,210 & \$13,000--\$13,500 & $>100\%$ \\
\bottomrule
\end{tabular}
\end{table}

\paragraph{Interpretation:}
\begin{itemize}
    \item \textbf{Discount Rate Variation:} Even at 10\%, NPW remains strongly positive, IRR above 90\%. The system still breaks even within about a year.
    \item \textbf{Labor Rate Variation:} Higher labor valuations yield proportionally higher NPW and IRR. If volunteer labor is $0/hr$, then water savings alone are insufficient to justify the project purely on a financial basis.
\end{itemize}

\subsection*{Conclusion: ``Good'' vs. ``Bad''}
All the tables and analyses confirm that when the reallocated volunteer labor is monetized (i.e., treat those hours as a real economic benefit), the project becomes extremely profitable: NPW surpasses \$12k, IRR exceeds 100\%, PI is over 10, and break-even happens inside the first year. However, absent any labor cost valuation, direct water savings of about \$30/year are too small to recoup \$1,330. 

Ultimately, from a broader sustainability and operational standpoint, the project is “good,” as it addresses multiple intangible benefits—environment, volunteer morale, educational opportunities—while also providing a strong “paper” ROI if volunteer hours are considered.

\newpage

\section*{Conclusions and Recommendations}
\vspace{-0.5em}

\subsection*{Analyses Performed}
\begin{itemize}
    \item Computed Break-even metrics (both simple and discounted) for a 10-year horizon.
    \item Calculated Net Present Worth (NPW) with the refined cost inputs.
    \item Determined the Internal Rate of Return (IRR), incorporating reallocated labor value and updated water-rate assumptions.
    \item Examined scenario-based variations in discount rate, labor costs, and component replacements (e.g., solar panel in year~6).
\end{itemize}

\subsection*{Breakeven Timeframe}
Even with the system’s initial capital outlay of approximately \$1,330 and a relatively small recurring cost (\$100/year), the system breaks even in under one year based on the reallocated labor value and modest water savings. Discounting at 7\% similarly indicates payback within or near the first year. 

\subsection*{Project Impact}
\begin{itemize}
    \item \textbf{High Net Present Worth (NPW):} 
    The NPW is approximately \$13,776 at a 7\% discount rate—showing a significant long-term benefit when incorporating volunteer labor as an effective cost offset.
    \item \textbf{Extraordinarily High IRR:} 
    Numerically, the IRR can exceed 100\% because the annual inflows (mainly from reallocated labor) are large relative to the modest up-front cost.
\end{itemize}
Although the raw financials are driven by the labor component, the results clearly indicate a strong economic rationale for implementing the new system.

\subsection*{Other Benefits}
\begin{itemize}
    \item \textbf{Sustainability \& Branding:} Demonstrates USC's leadership in water conservation, supporting its broader environmental stewardship goals.
    \item \textbf{Improved Reliability:} Replaces a prior design that failed after only a few months, ensuring consistent irrigation availability.
    \item \textbf{Educational Value:} Offers real-world lessons in sustainable infrastructure and system maintenance for students, volunteers, and staff.
    \item \textbf{Enhanced Volunteer Experience:} Frees up time previously spent on manual watering, enabling volunteers to focus on more engaging or impactful tasks.
\end{itemize}

\subsection*{Implementation Recommendations}
\begin{itemize}
    \item \textbf{Prompt Deployment:} Given the rapid payback and significant intangible benefits, installing the new system quickly maximizes overall returns.
    \item \textbf{Maintenance Planning:} Schedule annual checks of pump, battery (if used), and key mechanical connections; prepare a contingency fund of about \$200 for solar panel replacement in year~6--7.
    \item \textbf{Adapt for Future Expansion:} If additional campus gardens adopt similar designs, consider economies of scale when ordering parts or standardizing procedures.
    \item \textbf{Document Operating Procedures:} Provide straightforward winterization and debris-removal guidelines to ensure the system’s 10+ year service life.
\end{itemize}

\newpage
\section*{Appendix / References}
\begin{itemize}
    \item \textbf{Local Utility Rates:}  
    \href{http://www.columbiascwater.net}{City of Columbia Water} (effective July 1, 2024). Water rate used: \$0.0045 per gallon.
    \item \textbf{Labor Wages:}  
    Based on \href{https://www.bls.gov/oes/current/oes_sc.htm}{Bureau of Labor Statistics data for South Carolina} and regional job listings, sponsor wage structures typically range from \$16--\$18/hr; rate used: \$17/hr.
    \item \textbf{Material Pricing:}  
    Bill of Materials verified at major retailers: 
    \href{https://www.homedepot.com/}{Home Depot}, 
    \href{https://www.lowes.com/}{Lowe’s}, 
    \href{https://www.harborfreight.com/}{Harbor Freight}, and 
    \href{https://www.amazon.com/}{Amazon}.
    \item \textbf{Rainfall Data:}  
    \href{https://www.noaa.gov/}{NOAA Climate Normals} for Columbia, SC; approximately 48 inches/year average.
    \item \textbf{Discount Rate:}  
    A rate of 7\% was used based on common university capital planning guidelines (range 6--8\%). For a general overview of discount rates, see \href{https://www.investopedia.com/terms/d/discountrate.asp}{Investopedia's explanation of discount rates}.
\end{itemize}


\end{document}
